% Paper 6 (GRPO-Registry resource) standalone abstract
\begin{abstract}
The deltas that define GRPO variants---DAPO's asymmetric clip, dynamic sampling,
and token-level loss; Dr.~GRPO's dropped length and std normalization; GSPO's
sequence-level importance ratio---live in paper appendices and framework
defaults, not in any machine-readable form. The consequence, documented across
the \ours{} pillar papers, is that an algorithm \emph{label} underdetermines the
\emph{update rule} actually executed: in our audits, swapping the training backend
under a fixed label---which, undisclosed, also bundled a different base checkpoint,
so the label alone concealed \emph{both} changes---moved final reward across a
$17\times$ span between $85.6\%$ and $5.0\%$ (the same exhibit as the Pillar-5
backend audit, read here in the reverse direction),
and the same ``DAPO'' label yielded mean ZVF $0.00$ on an open trainer with true
dynamic sampling but $0.58$ on a closed stack running an asymmetric-clip
surrogate. These audits motivate \textsc{GRPO-Registry}, a resource that makes
such discrepancies queryable rather than presenting a new performance result:
(i) a JSON schema whose stack records carry the seven-field run-manifest block
of the eight-item MIN-REPORT-RL standard (loss form; reference policy and KL handling;
sampler/backend/precision \textbf{including base-checkpoint identity: revision/hash};
per-step ZVF/GU telemetry; group-size schedule;
held-out split; decontamination and parser probes), with held-out pass@$k$
curves recorded as the eighth, evaluation-report requirement, plus per-component
variant-delta status; (ii) three variant-delta records verified against the
DAPO, Dr.~GRPO, and GSPO papers; (iii) twenty seed stack entries derived from
this benchmark's released config dumps, an open Colab trainer audit, a
managed-runtime head-to-head, a same-stack four-method run, and a five-seed
risk-index batch; and (iv) a reference CLI implementing registry
queries, a 0--100 MIN-REPORT auditor badge, and \texttt{stackdiff}, which maps a
pair of entries to a flip-risk verdict (R0--R5). On the seed entries the badge
spans 43--96, and \texttt{stackdiff} flags the open-vs-closed ``DAPO'' pair as
an R5 label-flip risk from manifests alone. The registry, schema, and CLI are
released with the benchmark.
\paragraph{Venue claim boundary (position-artifact / resource).}
This is a \emph{machine-readable resource paper}, not a second independent
ranking of GRPO variants. The $17\times$ same-label exhibit is intentionally
shared with the companion MIN-REPORT-RL position paper (Pillar~5) and is cited
here only to motivate the schema; we do not re-claim it as novel empirical
science. The load-bearing contribution is the seven-field stack schema, seed
catalog, auditor badge, and \texttt{stackdiff} tooling. Submission target is a
position-artifact or datasets/resources track; living-catalog growth is
post-acceptance community work, not a gate on releasing the schema and seeds.
\paragraph{Absorbed satellite (retired).}
The former condensed living-catalog draft (roster R07) is retired from the live
tree to \texttt{archive/absorbed/R07\_grpo\_registry/}. Do not count R07 as a
separate venue paper.
\end{abstract}
\paragraph{August 2026 evidence boundary.} Registry entries report stated
provenance and conflicts; they do not reconcile unresolved exact-ID/model
conflicts or validate nominal algorithm comparisons.
